% ============================================================
% CHAPTER 7: CONCLUSION
% ============================================================
\chapter{Conclusion}
\label{ch:conclusion}

This dissertation set out to address the absence of a data-driven, network-level methodology for evaluating cycling infrastructure quality in England. While LTN~1/20 established the most comprehensive national design standards to date, its evaluation instruments were conceived for scheme-level audit and cannot assess whether a city's network, taken as a whole, delivers the connected low-stress conditions that its ``cycling for all'' principle requires. In response, this study developed a directed Level of Traffic Stress classification grounded entirely in LTN~1/20's own criteria and applied it to the West Midlands using the newly released OS NGD Transport data.

The contributions of this study span methodology and data. Methodologically, the classification replaces North American LTS thresholds with those derived from the CLoS Safety subcriteria, the JAT and the design criteria, producing the first LTS framework calibrated to English design standards and thereby restoring the population-anchored logic that LTN~1/20's compensatory scoring obscures. The directed network representation and bearing-based junction assessment further extend LTS practice beyond undirected max-propagation, incorporating roundabout-specific criteria absent from all prior implementations. On the data side, the study provides the first independent assessment of the OS NGD Cycle Lane dataset, establishing both its strengths in recording separated infrastructure and its systematic undercapture of paint-only provision, an evaluation that conditions the reliability of any future research built on this data source.

The findings demonstrate that abundant low-stress length does not constitute a usable network. Although low-stress segments dominate the West Midlands network by length, they fragment into ten thousands of disconnected components, with the largest contiguous low-stress network representing a small fraction of the whole. Arterial roads emerge as the primary source of severance, imposing high stress both along their links and at the junctions where low-stress residential clusters must cross them; this junction-level stress concentrates at the topological bottlenecks between clusters, while roundabouts function as near-absolute points of severance. The reach analysis makes the consequence tangible at the level of the individual cyclist: where on-carriageway provision lacks physical segregation, reachable distance without exceeding low-stress conditions collapses to that of the immediate residential cluster, terminating where the first arterial crossing is encountered. The result is that neither children nor mainstream adults can rely on the network for everyday journeys: only 4.1\% of the region's population can reach their nearest school via LTS~1 routes, and only 13.0\% of inter-MSOA commuting flow is connected via LTS~2 routes, with connected journeys requiring substantial detours.

These findings carry direct implications for investment practice. The growth simulation shows that the sequence of upgrades shapes connectivity outcomes as strongly as their total extent, with a strategy targeting component-bridging edges reaching connectivity milestones in roughly half the upgrade length of uninformed selection. This exposes a gap in the current policy architecture: LTN~1/20's funding-eligibility tools evaluate individual schemes, yet the connectivity of the network depends on where and in what order schemes are delivered, a coordination problem that scheme-level appraisal cannot resolve.

Because both the classification thresholds and the underlying data are national in scope, the framework can be applied to any English city without recalibration, enabling systematic benchmarking of network conditions against a single statutory standard. Future work should validate the classification against cyclist-perceived stress, extend junction assessment through movement-level dual-graph representations, and incorporate cost and equity dimensions into growth modelling. Ultimately, this dissertation demonstrates that LTN~1/20 contains within itself the material for its own network-scale evaluation. By translating its design criteria into a stress classification, the standard can be turned back upon the networks it governs, revealing in quantitative terms the distance between the ambition of cycling for all and the infrastructure through which that ambition must be realised.
